\chapter{Dermatomyositis}
\index{Dermatomyositis}
\label{sec:skin:Dermatomyositis}

\section{Overview}
Dermatomyositis is an idiopathic inflammatory myopathy with characteristic skin manifestations, affecting both children and adults.
\section{Causes}
This condition happens because of autoimmune-mediated inflammation of skin and muscle. Cutaneous features include heliotrope rash (violaceous discoloration of the eyelids), Gottron's papules (erythematous to violaceous papules overlying the knuckles), V-sign (erythematous rash on the anterior neck and chest), shawl sign (erythematous rash over shoulders and upper back), mechanic's hands (rough, cracked skin on the fingers), and nail fold capillary changes. Proximal muscle weakness is the hallmark of myopathy.     
\section{Risk Factors}

\begin{itemize}[noitemsep]
  \item Genetic predisposition and family history
  \item Advancing age
  \item Lifestyle factors (diet, exercise, smoking, alcohol)
  \item Environmental exposures
  \item Underlying medical conditions
  \item Medication use
\end{itemize}
\section{Signs and Symptoms}

\subsection{Common symptoms}
\begin{itemize}[noitemsep]
  \item You may experience both skin and muscle manifestations
  \item Dermatomyositis in adults carries an increased risk of underlying malignancy (ovarian, lung, GI), necessitating age-appropriate cancer screening.
  \item Pain or discomfort in the affected area
  \end{itemize}

\subsection{Emergency warning signs}
\begin{itemize}[noitemsep]
  \item Severe or worsening symptoms of dermatomyositis require immediate medical evaluation
\end{itemize}
\section{Progression and Stages}
The condition may progress through different stages of severity over time. Early stages often involve mild or subtle symptoms that may be overlooked. Moderate stages involve more noticeable symptoms affecting daily function. Advanced stages may involve significant impairment and complications.
\section{Complications}
\begin{itemize}[noitemsep]
  \item Disease progression leading to worsening symptoms
  \item Secondary infections or injuries
  \item Side effects of treatment
  \item Reduced quality of life
  \item Emotional and psychological impact
\end{itemize}
\section{Diagnosis}
Diagnosis typically involves CK levels, myositis-specific antibodies (anti-Mi-2, anti-MDA5, anti-TIF1-$\gamma$), muscle MRI, EMG, and muscle biopsy. Comprehensive medical history and physical examination Laboratory tests to assess organ function and rule out other conditions Imaging studies to visualize affected structures Specialized testing depending on the affected system Consultation with relevant specialists Monitoring of disease progression over time

\begin{itemize}[noitemsep]
  \item Comprehensive medical history and physical examination
  \item Laboratory tests to assess organ function
  \item Imaging studies to visualize affected structures
  \item Specialized testing depending on the affected system
  \item Consultation with relevant specialists
\end{itemize}
\section{Prevention}
\begin{itemize}[noitemsep]
  \item Maintain a healthy lifestyle with balanced diet and exercise
  \item Avoid known risk factors
  \item Attend regular medical check-ups for early detection
  \item Follow recommended screening guidelines
\end{itemize}
\section{Treatment Options}
Treatment options include corticosteroids, steroid-sparing agents (methotrexate, azathioprine, mycophenolate), IVIG, and rituximab. Medications to manage symptoms and address underlying mechanisms Lifestyle modifications and self-care measures Physical therapy and rehabilitation Surgical intervention when indicated Regular monitoring and follow-up care Patient education and support services

\begin{itemize}[noitemsep]
  \item Medications to manage symptoms and address underlying mechanisms
  \item Lifestyle modifications and self-care measures
  \item Physical therapy and rehabilitation
  \item Surgical intervention when indicated
  \item Regular monitoring and follow-up care
\end{itemize}
\section{Living with It}
\begin{itemize}[noitemsep]
  \item Follow your treatment plan as prescribed
  \item Maintain regular follow-up with your healthcare provider
  \item Make necessary lifestyle adjustments
  \item Seek support from family, friends, or support groups
  \item Stay informed about your condition
\end{itemize}
\section{Outlook}
\begin{itemize}[noitemsep]
  \item The outlook varies from person to person
  \item Malignancy screening is essential in adults.
\end{itemize}
